\documentclass[11pt]{article}
\usepackage[a4paper,margin=26mm]{geometry}
\usepackage{amsmath,amssymb,mathtools}
\usepackage{unicode-math}
\setmainfont{Latin Modern Roman}
\setmathfont{Latin Modern Math}
\usepackage{xurl,hyperref}
\hypersetup{hidelinks,pdftitle={UnsolvedMath Problem 10300014}}
\newcommand{\sref}[2]{\hyperlink{source:#1:#2}{[S#2]}}
\newcommand{\cataloguescope}[1]{#1}
\begin{document}
% UnsolvedMath ID 10300014; source catalogue code AMR-102-0014
\setcounter{section}{10300013}
\section{Sublaminations and superlaminations — Question 6.5}\label{10300014}
{\small\sffamily UnsolvedMath ID 10300014 · Topology}
\subsection{Definitions and mathematical statement}

\paragraph{Shared notation.}$\mathbb N=\{1,2,\ldots\}$, and $\mathbb Z,\mathbb Q,\mathbb R,\mathbb C$ denote the integers, rationals, reals and complex numbers. Any other conventions are specified locally.


A surface lamination $\Lambda$ is a closed subset of $M$ locally homeomorphic to a product of a disk with a closed transverse set. It is essential when it has no sphere leaf or torus leaf bounding a solid torus, and its complementary regions are irreducible, have incompressible boundary, and have no compressing monogon. Such a monogon is a compressing disk with one boundary arc on a boundary leaf and the other along a transverse interval fiber. It is genuine if some complementary region is not an interval bundle. A sublamination is a closed union of leaves. Write $\widetilde M$ for the universal cover of $M$ and $\widetilde\Lambda$ for the full inverse-image lamination under the covering map. A loosesse lamination obeys the essential-lamination conditions except that compressing monogons are allowed. A leaf is properly embedded if its inclusion is an embedding and the inverse image of every compact ambient set is compact.
\paragraph{Statement recorded in the supplied source.}
Determine useful topological consequences of a loosesse lamination. In particular, must every leaf of $\widetilde\Lambda$ be properly embedded in $\widetilde M$? Does existence of such a lamination imply $\widetilde M\cong\mathbb R^3$? The broader utility question is part of the task, in addition to these two explicit tests. \sref{10300014}{2}
\subsection{Short English statement}
Determine the usefulness of loosesse laminations, whether their lifted leaves are proper, and whether they force the universal cover to be Euclidean three-space.
\subsection{Sources}
\begin{description}
\item[\hypertarget{source:10300014:1}{S1}] ULAM.AI and the UnsolvedMath Contributors, \emph{UnsolvedMath: A Curated Collection of Open Mathematics Problems}, record 10300014 (AMR-102-0014). CC BY 4.0. \url{https://huggingface.co/datasets/ulamai/UnsolvedMath}.
\item[\hypertarget{source:10300014:2}{S2}] Danny Calegari, \emph{Problems in foliations and laminations of 3-manifolds}, version 0.78 (2002), Question 6.5; Definitions 1.1--1.10 and local definitions and remarks. \url{https://arxiv.org/abs/math/0209081v1}.
\end{description}
\label{end:10300014}
\end{document}
